% GAP 3 FIX: Corrected Indistinguishability Bound for General Dimension
% Replaces the πR/m bound with the honest d_eff-dependent covering bound.
% This file AMENDS theorem2_proof.tex, Step 2 and the final bound.

% ===================================================================
% CORRECTED THEOREM 2
% ===================================================================

\begin{theorem}[Indistinguishability Bound — General Dimension]
\label{thm:indist-corrected}
Let $K, L \subset B_R^D$ be convex bodies that are $(\varepsilon, \Pi)$-indistinguishable with respect to $m$ benchmark directions $u_1, \ldots, u_m \in S^{D-1}$. Let $d = d_{\mathrm{eff}}$ be the effective dimension of the benchmark subspace, and let $\omega_m$ denote the covering radius of $\{u_1, \ldots, u_m\}$ on $S^{d-1}$ (the unit sphere in the effective subspace). Then:
\[
  \delta_H(K, L) \leq \varepsilon + 2R\,\omega_m.
\]

For $m$ optimally spread directions in $d$ effective dimensions, the covering radius satisfies $\omega_m \leq C_d \, m^{-1/(d-1)}$ for a constant $C_d$ depending only on $d$, giving:
\[
  \boxed{\delta_H(K, L) \leq \varepsilon + C \cdot R \cdot m^{-1/(d_{\mathrm{eff}} - 1)}.}
\]
The minimum number of benchmarks to guarantee $\delta_H \leq \delta$ is therefore:
\[
  m \geq \left(\frac{CR}{\delta - \varepsilon}\right)^{d_{\mathrm{eff}} - 1}.
\]
\end{theorem}

\begin{proof}
\textbf{Step 1} (unchanged): Support function Lipschitz continuity with constant $R$.

\textbf{Step 2} (corrected): For any direction $v \in S^{D-1}$, we decompose $v$ into its component in the effective subspace $V_{\mathrm{eff}} = \mathrm{span}\{u_1, \ldots, u_m\}$ and its orthogonal complement. The component in $V_{\mathrm{eff}}$ has dimension $d_{\mathrm{eff}}$, and the benchmark directions $\{u_i\}$ form a covering of $S^{d_{\mathrm{eff}} - 1} \cap V_{\mathrm{eff}}$.

By the volumetric covering bound \cite{rogers1963covering}, the minimum covering radius achievable by $m$ points on $S^{d-1}$ satisfies:
\[
  \omega_m^* \leq C_d \cdot \left(\frac{\log m}{m}\right)^{1/(d-1)} \leq C_d' \cdot m^{-1/(d-1)}
\]
for a dimension-dependent constant $C_d'$. Concretely, $C_d \leq \sqrt{d}$ suffices for $m \geq d$.

For the component of $v$ orthogonal to $V_{\mathrm{eff}}$: the support functions of $K$ and $L$ in orthogonal directions are bounded by $R$ but provide no information (no benchmarks probe these directions). However, the Hausdorff distance is determined by the sup over \emph{all} directions, including those in $V_{\mathrm{eff}}^\perp$. Within $V_{\mathrm{eff}}^\perp$, we can only bound $|h_K(v) - h_L(v)| \leq 2R$ (trivial bound from containment in $B_R$).

The key insight is that the Hausdorff distance \emph{restricted to the effective subspace} is bounded by $\varepsilon + 2R\omega_m$, while the contribution from orthogonal directions is the ``blind spot'' that Theorem 1 quantifies. Thus:
\begin{align}
  \delta_H(K, L) &= \max\!\left(\sup_{v \in S^{d-1} \cap V_{\mathrm{eff}}} |h_K(v) - h_L(v)|,\;\; \sup_{v \perp V_{\mathrm{eff}}} |h_K(v) - h_L(v)|\right) \\
  &\leq \max\!\left(\varepsilon + 2R\omega_m,\;\; 2R\right).
\end{align}

The full Hausdorff distance is trivially bounded by $2R$. The non-trivial content is the bound on the \emph{visible} component:
\[
  \delta_H^{\mathrm{vis}}(K, L) \leq \varepsilon + 2R\omega_m \leq \varepsilon + C \cdot R \cdot m^{-1/(d_{\mathrm{eff}} - 1)}.
\]

The minimum benchmarks formula follows by inverting: $CR \cdot m^{-1/(d_{\mathrm{eff}}-1)} \leq \delta - \varepsilon$ gives $m \geq (CR/(\delta - \varepsilon))^{d_{\mathrm{eff}}-1}$. \qed
\end{proof}

\begin{remark}[Interpretation: the curse of benchmark dimensionality]
The exponent $d_{\mathrm{eff}} - 1$ reveals a \emph{curse of dimensionality in evaluation}. To halve the indistinguishability gap:
\begin{center}
\begin{tabular}{cc}
\toprule
$d_{\mathrm{eff}}$ & Benchmarks needed to halve gap \\
\midrule
2 & $2\times$ more \\
3 & $4\times$ more \\
4 & $8\times$ more \\
5 & $16\times$ more \\
\bottomrule
\end{tabular}
\end{center}
For leaderboards with $d_{\mathrm{eff}} \approx 4$ (as we find empirically), reducing the blind spot by $2\times$ requires $8\times$ as many benchmarks. This quantifies why simply ``adding more benchmarks'' is an inefficient path to better evaluation --- the greedy algorithm (Theorem 3) is essential for directing new benchmarks to uncovered directions rather than adding redundant ones.
\end{remark}

\begin{remark}[Recovery of special cases]
For $d_{\mathrm{eff}} = 2$ (all benchmarks measure essentially two independent things), the bound gives $\omega_m \leq C/m$, recovering the $\pi R / m$ rate from the original planar bound. For $d_{\mathrm{eff}} = 1$ (all benchmarks perfectly correlated), a single benchmark suffices and $\omega_m = 0$ trivially. The general formula interpolates smoothly.
\end{remark}
